% GENERATED FILE. Edit canonical lessons, then run npm run generate:books.
% canonical-source-hash: fnv1a64:42826416c9c038bd
% canonical-lessons: ML-C254-kuruvi, ML-C254-annan, ML-C254-nantu, ML-C254-val, ML-C254-ciraku

\chapter{Sparrow, Squirrel, Crab, Tail, Wing}
\label{ch:ml-sparrow-squirrel-crab-tail-wing}
\hlchaptermodality{254}

\begin{chapteropening}
\textbf{By the end of this chapter:} \emph{I can say a sparrow, a squirrel, a crab, a tail and a wing.}

Complete the last lesson of chapter 254: I can say a sparrow, a squirrel, a crab, a tail and a wing.
\end{chapteropening}

\section[\texorpdfstring{\ml{കുരുവി} (kuruvi)}{കുരുവി (kuruvi)}]{\ml{കുരുവി} (kuruvi) — a sparrow}

\label{lesson:ML-C254-kuruvi}



\begin{quote}
Before the new one: say the Malayalam for a hibiscus, then the Malayalam for holy basil.
\end{quote}

\subsection*{You'll want to know: \ml{കുരുവി}}
\textbf{\ml{കുരുവി}} — \emph{kuruvi} — \textquotedblleft{}a sparrow\textquotedblright{}.

\begin{cousinweb}
\textbf{\ml{കുരുവി}} (\emph{kuruvi}) is native Dravidian, one word with Tamil \emph{kuruvi} and Kannada \emph{gubbi} (DEDR 1793).
\end{cousinweb}

\begin{grammarlens}[title={What you've built}]
One more word for this chapter.
\end{grammarlens}

\subsection*{Your turn}
\begin{itemize}
\raggedright
  \item \emph{Say it:} \emph{kuruvi}
  \item \emph{Say it:} \emph{kuruvi}, once more
  \item \emph{Say it:} say \emph{kuruvi}
  \item \emph{Recall:} say the Malayalam for a hibiscus, then the Malayalam for holy basil, then say \emph{kuruvi} again
\end{itemize}

\begin{tcolorbox}[breakable,colback=teal!4,colframe=teal!35!black,title={Before you move on}]
What does \ml{കുരുവി} mean? (\textquotedblleft{}A sparrow\textquotedblright{}.) Say it once more.
\end{tcolorbox}
\section[\texorpdfstring{\ml{അണ്ണാൻ} (aṇṇān)}{അണ്ണാൻ (aṇṇān)}]{\ml{അണ്ണാൻ} (aṇṇān) — a squirrel}

\label{lesson:ML-C254-annan}



\begin{quote}
Before the new one: say the Malayalam for holy basil, then the Malayalam for a sparrow.
\end{quote}

\subsection*{You'll want to know: \ml{അണ്ണാൻ}}
\textbf{\ml{അണ്ണാൻ}} — \emph{aṇṇān} — \textquotedblleft{}a squirrel\textquotedblright{}.

\begin{cousinweb}
\textbf{\ml{അണ്ണാൻ}} (\emph{aṇṇān}) is native Dravidian, one word with Tamil \emph{aṇil} and Kannada \emph{aṇal} (DEDR 2315).
\end{cousinweb}

\begin{grammarlens}[title={What you've built}]
One more word for this chapter.
\end{grammarlens}

\subsection*{Your turn}
\begin{itemize}
\raggedright
  \item \emph{Say it:} \emph{aṇṇān}
  \item \emph{Say it:} \emph{aṇṇān}, once more
  \item \emph{Say it:} say \emph{aṇṇān}
  \item \emph{Recall:} say the Malayalam for holy basil, then the Malayalam for a sparrow, then say \emph{aṇṇān} again
\end{itemize}

\begin{tcolorbox}[breakable,colback=teal!4,colframe=teal!35!black,title={Before you move on}]
What does \ml{അണ്ണാൻ} mean? (\textquotedblleft{}A squirrel\textquotedblright{}.) Say it once more.
\end{tcolorbox}
\section[\texorpdfstring{\ml{ഞണ്ട്} (ñaṇṭŭ)}{ഞണ്ട് (ñaṇṭŭ)}]{\ml{ഞണ്ട്} (ñaṇṭŭ) — a crab}

\label{lesson:ML-C254-nantu}



\begin{quote}
Before the new one: say the Malayalam for a sparrow, then the Malayalam for a squirrel.
\end{quote}

\subsection*{You'll want to know: \ml{ഞണ്ട്}}
\textbf{\ml{ഞണ്ട്}} — \emph{ñaṇṭŭ} — \textquotedblleft{}a crab\textquotedblright{}.

\begin{cousinweb}
\textbf{\ml{ഞണ്ട്}} (\emph{ñaṇṭŭ}) is native Dravidian, one word with Tamil \emph{ñaṇṭu}, Kannada \emph{eṇḍrakāyi} and Telugu \emph{eṇḍri} (DEDR 2901).
\end{cousinweb}

\begin{grammarlens}[title={What you've built}]
One more word for this chapter.
\end{grammarlens}

\subsection*{Your turn}
\begin{itemize}
\raggedright
  \item \emph{Say it:} \emph{ñaṇṭŭ}
  \item \emph{Say it:} \emph{ñaṇṭŭ}, once more
  \item \emph{Say it:} say \emph{ñaṇṭŭ}
  \item \emph{Recall:} say the Malayalam for a sparrow, then the Malayalam for a squirrel, then say \emph{ñaṇṭŭ} again
\end{itemize}

\begin{tcolorbox}[breakable,colback=teal!4,colframe=teal!35!black,title={Before you move on}]
What does \ml{ഞണ്ട്} mean? (\textquotedblleft{}A crab\textquotedblright{}.) Say it once more.
\end{tcolorbox}
\section[\texorpdfstring{\ml{വാൽ} (vāl)}{വാൽ (vāl)}]{\ml{വാൽ} (vāl) — a tail}

\label{lesson:ML-C254-val}



\begin{quote}
Before the new one: say the Malayalam for a squirrel, then the Malayalam for a crab.
\end{quote}

\subsection*{You'll want to know: \ml{വാൽ}}
\textbf{\ml{വാൽ}} — \emph{vāl} — \textquotedblleft{}a tail\textquotedblright{}.

\begin{cousinweb}
\textbf{\ml{വാൽ}} (\emph{vāl}) $\leftarrow$ Sanskrit \emph{vāla}, \textquotedblleft{}tail, hair of the tail\textquotedblright{} (CDIAL 11572). DEDR appendix 57, which gathers Dravidian words taken from Indo-Aryan, lists it with Tamil \emph{vāl} and Kannada \emph{bāla}.
\end{cousinweb}

\begin{grammarlens}[title={What you've built}]
One more word for this chapter.
\end{grammarlens}

\subsection*{Your turn}
\begin{itemize}
\raggedright
  \item \emph{Say it:} \emph{vāl}
  \item \emph{Say it:} \emph{vāl}, once more
  \item \emph{Say it:} say \emph{vāl}
  \item \emph{Recall:} say the Malayalam for a squirrel, then the Malayalam for a crab, then say \emph{vāl} again
\end{itemize}

\begin{tcolorbox}[breakable,colback=teal!4,colframe=teal!35!black,title={Before you move on}]
What does \ml{വാൽ} mean? (\textquotedblleft{}A tail\textquotedblright{}.) Say it once more.
\end{tcolorbox}
\section[\texorpdfstring{\ml{ചിറക്} (ciṟakŭ)}{ചിറക് (ciṟakŭ)}]{\ml{ചിറക്} (ciṟakŭ) — a wing}

\label{lesson:ML-C254-ciraku}



\begin{quote}
Before the new one: say the Malayalam for a crab, then the Malayalam for a tail.
\end{quote}

\subsection*{You'll want to know: \ml{ചിറക്}}
\textbf{\ml{ചിറക്}} — \emph{ciṟakŭ} — \textquotedblleft{}a wing\textquotedblright{}.

\begin{cousinweb}
\textbf{\ml{ചിറക്}} (\emph{ciṟakŭ}) is native Dravidian, one word with Tamil \emph{ciṟai}, Kannada \emph{eṟake} and Telugu \emph{eṟaka} (DEDR 2591).
\end{cousinweb}

\begin{grammarlens}[title={What you've built}]
One more word for this chapter.
\end{grammarlens}

\subsection*{Your turn}
\begin{itemize}
\raggedright
  \item \emph{Say it:} \emph{ciṟakŭ}
  \item \emph{Say it:} \emph{ciṟakŭ}, once more
  \item \emph{Say it:} say \emph{ciṟakŭ}
  \item \emph{Recall:} say the Malayalam for a crab, then the Malayalam for a tail, then say \emph{ciṟakŭ} again
\end{itemize}

\begin{tcolorbox}[breakable,colback=teal!4,colframe=teal!35!black,title={Before you move on}]
What does \ml{ചിറക്} mean? (\textquotedblleft{}A wing\textquotedblright{}.) Say it once more.
\end{tcolorbox}
